\documentclass{ieeeaccess}
\usepackage{cite}
\usepackage{amsmath,amssymb,amsfonts}
\usepackage{algorithmic}
\usepackage{graphicx}
\usepackage{textcomp}
\usepackage{caption}
\usepackage{mathptmx}
\usepackage{hyperref}
\usepackage[english]{babel}
\usepackage{microtype} 
\usepackage{xurl}
\usepackage{booktabs}
\usepackage{tabularx}
\usepackage{placeins}
\usepackage{bm}
\usepackage{bbm}
\usepackage{tabularx}
\usepackage{ragged2e}
\makeatletter
\AtBeginDocument{\DeclareMathVersion{bold}
\SetSymbolFont{operators}{bold}{T1}{times}{b}{n}
\SetSymbolFont{NewLetters}{bold}{T1}{times}{b}{it}
\SetMathAlphabet{\mathrm}{bold}{T1}{times}{b}{n}
\SetMathAlphabet{\mathit}{bold}{T1}{times}{b}{it}
\SetMathAlphabet{\mathbf}{bold}{T1}{times}{b}{n}
\SetMathAlphabet{\mathtt}{bold}{OT1}{pcr}{b}{n}
\SetSymbolFont{symbols}{bold}{OMS}{cmsy}{b}{n}
\renewcommand\boldmath{\@nomath\boldmath\mathversion{bold}}}
\makeatother

\def\BibTeX{{\rm B\kern-.05em{\sc i\kern-.025em b}\kern-.08em
    T\kern-.1667em\lower.7ex\hbox{E}\kern-.125emX}}

%Your document starts from here ___________________________________________________
\begin{document}
\history{Date of publication xxxx 00, 0000, date of current version xxxx 00, 0000.}
\doi{10.1109/ACCESS.2024.0429000}


\title{On Reliable Detection of AI-Generated Multimedia Content}
\author{\uppercase{Ansuman Nayak}\authorrefmark{1}, and 
\uppercase{Vijay K. Madisetti}\authorrefmark{2}, \IEEEmembership{Fellow, IEEE}}

\address[1]{College of Computing, Georgia Institute of Technology, Atlanta, GA 30332 USA (e-mail: anayak306@gatech.edu)}
\address[2]{School of Cybersecurity and Privacy, Georgia Institute of Technology, Atlanta, GA 30332 USA (e-mail: vkm@gatech.edu)}

\markboth
{Nayak \headeretal: On Reliable Detection of AI-Generated Multimedia Content}
{Nayak \headeretal: On Reliable Detection of AI-Generated Multimedia Content}

\corresp{Corresponding author: Dr. Vijay Madisetti (e-mail: vkm@gatech.edu).}


\begin{abstract}
With the rapid proliferation of AI-generated multimedia content, the growing need for reliable detection of AI-generated media, including deepfakes and misinformation, has been largely unmet. Uniformly applying large vision-language models (VLMs) to all content can be prohibitively expensive. In this work, we formulate detection as a model selection problem over a set of complementary detectors. We show that disagreement between artifact-based and semantic detectors provides a strong, parameter-free indicator of input difficulty, and that this signal remains stable even when the semantic probe faces deployment shift to unseen generator families. We construct a selective-computation cascade in which lightweight detectors process all inputs and a stronger verifier is invoked only for disagreement cases. On in-domain benchmarks, a well-trained semantic detector achieves the best cost-accuracy tradeoff without VLM escalation. The cascade's utility emerges under generator evolution, but only when the verifier remains stronger than the stale semantic detector on the routed disagreement subset. On SDXL and FLUX, escalating to GPT-5.5 only on disagreement cases recovers up to 4.9 percentage points of accuracy over a stale semantic probe while avoiding 58--65\% of verifier invocations. In contrast, under a more distant BigGAN shift, escalation provides no benefit, showing that routing effectiveness depends on verifier strength and distributional distance rather than disagreement alone. We further introduce a verifier validation gate that correctly predicts whether escalation is beneficial in 20 of 21 detector-verifier configurations before deployment.

\end{abstract}

\begin{keywords}
AI-generated image detection, deepfake 
detection, multi-model disagreement, distribution shift, 
vision-language models, forensic routing, CLIP embeddings
\end{keywords}

\titlepgskip=-21pt
\maketitle

% Sections
\section{Introduction}

The rapid advancement of generative models for images and videos has significantly increased the realism and accessibility of synthetic media. Diffusion models~\cite{Wang_2023_ICCV}, GAN-based generators~\cite{wang2020cnn}, and multimodal large language models (MLLMs) can now generate highly convincing content with minimal effort. Although this opens up useful applications, it also raises concerns about misinformation, impersonation, and fraud. As a result, reliably detecting AI-generated media has become an important problem that still needs to be addressed.

One of the main challenges is generalization. Detection models often perform well on the generator families they were trained on, but accuracy degrades significantly when applied to new or unseen generators~\cite{NEURIPS2023_f4d4a021}. Artifact-based detectors exploit low-level frequency and synthesis patterns introduced during generation, while semantic models capture higher-level structural representations. These two signal types fail in different ways, suggesting that combining them provides more robust coverage than either approach alone.


In practical deployment settings, compute cost is a significant concern. High-capacity
vision-language models (VLMs) provide strong reasoning capabilities but are expensive to apply at scale~\cite{chen2024frugalgpt}. Applying these models uniformly to every input is often infeasible in real-time or large-scale systems. Existing approaches typically apply a
fixed detector or ensemble without adapting computation to input difficulty. The central challenge is identifying which inputs require expensive verification rather than applying a
uniform computation to all inputs.


In this work, we address this challenge by modeling AI-generated media detection as a model selection problem over complementary detectors. We find that the disagreement between an artifact-based detector and a semantic detector serves as a proxy for input difficulty and requires no calibration. When the detectors agree, the predictions are highly reliable, whereas disagreements consistently identify cases where cheap detectors are insufficient.

Unlike prior approaches that rely on calibrated confidence scores or learned routing policies, disagreement provides a parameter-free signal that remains reliable under distribution shift. We design a cascade that routes inputs based on a signal where lightweight detectors process all inputs, and a stronger verifier is invoked only on disagreement cases. When the semantic probe is trained in-domain, it already achieves strong performance without VLM escalation. The cascade's primary value is under deployment shift, where it recovers
accuracy over a stale semantic probe while avoiding the majority of verifier invocations. This reframes the detection problem as identifying which inputs require expensive computation, rather than uniformly increasing model complexity.

\subsection{Research Questions}

This work is guided by the following research questions:

\begin{itemize}
    \item \textbf{RQ1:} Can we reliably estimate which inputs require more powerful models
    for AI-generated media detection, including under generator distribution shift?
    \item \textbf{RQ2:} Can such estimates be used to guide model selection across heterogeneous detectors without relying on calibrated confidence scores?
    \item \textbf{RQ3:} Does disagreement-based routing reduce verifier cost while recovering accuracy under deployment shift to unseen generator families?
\end{itemize}

\subsection{Contributions}

This work makes the following contributions:

\begin{itemize}
    \item \textbf{C1:} We show that the disagreement between artifact-based and semantic detectors reliably predicts failure cases across generator families and remains
    predictive under deployment shift to unseen generators.
    
    \item \textbf{C2:} We propose a parameter-free routing strategy that uses disagreement as a difficulty signal for selective model invocation, requiring no confidence
    calibration or policy training.
    
    \item \textbf{C3:} On deployment shift benchmarks (SDXL and FLUX), cascade escalation to GPT-5.5 recovers up to 4.9 percentage points of accuracy over a stale semantic probe while reducing verifier invocations by 58--65\%.
    
    \item \textbf{C4:} We show that a learned contextual bandit converges to the disagreement rule across multiple datasets, independently validating disagreement as the dominant routing signal while providing no measurable improvement over the parameter-free rule.

    \item \textbf{C5:} We introduce a verifier validation gate that determines whether disagreement-based escalation should be enabled before deployment. Across 21 detector-verifier configurations, including a distant GAN-to-diffusion shift (BigGAN), it correctly predicts whether escalation improves accuracy in 20 cases (95.2\%), providing a practical deployment safeguard.
\end{itemize}


\section{Related Work}
This section reviews prior work related to AI-generated media detection and adaptive inference. We first identify gaps in existing approaches that motivate this work, then discuss detection methods under distribution shift, examine approaches for selective computation, and finally compare routing strategies based on confidence and learned policies.




\subsection{Gaps in Existing Approaches}

Existing forensic detection and cost-aware inference literature leaves several gaps that motivate this work.

\textbf{Single signal reliance.} Most forensic detectors rely on a single signal type, such as diffusion-based artifacts or reconstruction inconsistencies~\cite{Corvi_2023_ICASSP} or semantic embeddings from large pretrained models~\cite{cozzolino2023raising}. These approaches work well within specific generator families but degrade under generator distribution shift. Recent large-scale evaluations show that no single detector consistently generalizes across modern image generation pipelines~\cite{epstein2023online, cozzolino2023raising}.

\textbf{Confidence-based routing under shift.} Cost-aware inference methods and selective classification approaches rely on calibrated confidence estimates to decide when to defer or escalate predictions. Predictive uncertainty in deep models degrades significantly under distribution shift~\cite{ovadia2019can}. As a result, confidence-based routing signals become unreliable in the very regimes where robust detection is most needed. Existing literature on uncertainty estimation and selective prediction does not address this failure mode in heterogeneous detector cascades.

\textbf{Disagreement as a routing signal.} Inter-model disagreement has been studied extensively in active learning~\cite{seung1992query} and ensemble-based uncertainty estimation~\cite{lakshminarayanan2017simple}. In these settings, disagreement identifies informative samples or improves predictive uncertainty. Prior work treats disagreement as a mechanism for aggregation or data selection. Its role as a routing signal for cost-aware inference remains unexplored in forensic detection pipelines. Specifically, using disagreement as a diagnostic to indicate when low-cost detectors should not be trusted has not been studied.

This work addresses these gaps in four ways. First, we combine heterogeneous artifact-based and semantic signals through a parameter-free disagreement rule. Second, we demonstrate empirically that disagreement provides a more reliable routing signal than confidence under distribution shift. We evaluate this at matched operating points across multiple generator families. Third, we account for both latency and token cost when evaluating VLM-based escalation strategies. Fourth, we reinterpret disagreement as a diagnostic for selective invocation of more expensive models, not as a voting mechanism.

\subsection{Detection under Distribution Shift}

A large body of work has focused on detecting AI-generated images using artifact-based features. Early approaches such as CNN-based detectors~\cite{wang2020cnn} exploit low-level patterns introduced by generative models. More recent methods aim to improve generalization across generators, such as universal fake detectors~\cite{ojha2023fakedetect}. Benchmark datasets such as GenImage~\cite{NEURIPS2023_f4d4a021} highlight the difficulty of this problem, showing that many detectors degrade significantly when evaluated on unseen generators.

Other approaches leverage higher-level representations. Models such as CLIP~\cite{pmlr-v139-radford21a} learn semantic features from large-scale image-text data and have been used in downstream detection tasks. Diffusion-specific detection methods, such as DIRE~\cite{Wang_2023_ICCV}, further illustrate the diversity of signals used in modern detection systems. Despite these advances, no single model consistently generalizes across generator families, motivating the use of multiple complementary detectors.

Frequency-domain analysis provides a complementary artifact signal. Spectral methods
exploit periodic and structured artifacts introduced by GAN upsampling and diffusion
denoising pipelines, and have shown strong in-domain generalization~\cite{Corvi_2023_ICASSP}.
Beyond CLIP, recent vision encoders such as DINOv2~\cite{oquab2023dinov2} and
SigLIP2~\cite{tschannen2025siglip2} offer stronger visual representations and have
demonstrated improved robustness under limited training data. Their behavior under
generator distribution shift in forensic detection tasks remains underexplored, motivating
their inclusion as semantic detector candidates in this work.


\subsection{Selective Computation and Model Cascades}

Selective computation has been studied as a way to reduce inference cost by adapting computation to input difficulty. Early work on selective classification~\cite{NIPS2017_4a8423d5} proposes abstaining on uncertain predictions to improve reliability. More recent systems extend this idea to model selection, routing inputs across models of varying capacity. For example, FrugalGPT~\cite{chen2024frugalgpt} reduces cost by selecting among large language models based on estimated utility. These approaches typically rely on confidence scores or calibrated probabilities to guide routing decisions. However, such signals can be unreliable under distribution shift, where models may produce overconfident incorrect predictions.

\subsection{Routing Signals for Model Selection}

Most existing routing strategies use confidence-based signals, such as entropy or prediction probability, to determine whether additional computation is required. While effective in controlled settings, these signals often degrade under distribution shift due to poor calibration. Learned routing policies, such as bandit-based approaches, attempt to improve decision-making by training a model to predict the best action
from input features.

In contrast, this work explores disagreement between heterogeneous detectors as a routing signal. Rather than relying on calibrated confidence scores or learned policies, we use prediction disagreement between artifact-based and semantic models to estimate input difficulty. This provides a simple, calibration-free signal that remains effective under distribution shift.

\begin{table*}[h]
\centering
\small % Reduces font size to help fit more content comfortably
\caption{Comparative analysis of related AI-generated media detection approaches. \checkmark~indicates capability is present; \texttimes~indicates absent.}
\label{tab:litsurvey}
\setlength{\tabcolsep}{4pt} % Slightly tighter spacing between columns

% Define a custom column type 'Y' that is a ragged-right version of 'X'
\newcolumntype{Y}{>{\RaggedRight\arraybackslash}X}

\begin{tabularx}{\textwidth}{p{2.0cm}YYYYYY}
\toprule
\textbf{Feature}
  & \textbf{CNNSpot} \cite{wang2020cnn}
  & \textbf{DIRE} \cite{Wang_2023_ICCV}
  & \textbf{UnivFD} \cite{ojha2023fakedetect}
  & \textbf{BusterX} \cite{wen2025busterxmllmpoweredaigeneratedvideo}
  & \textbf{FrugalGPT} \cite{chen2024frugalgpt}
  & \textbf{Ours} \\
\midrule
Detection signal
  & (\checkmark) Frequency \& upsampling artifacts
  & (\checkmark) Diffusion reconstruction error
  & (\checkmark) CLIP embeddings for generalizable detection
  & (\checkmark) MLLM open-ended visual reasoning
  & (\texttimes) LLM text output, not a forgery detector
  & (\checkmark) Artifact + semantic embeddings + VLM reasoning \\
\midrule
Model architecture
  & (\texttimes) Single CNN/fixed computation
  & (\texttimes) Single diffusion model/fixed computation
  & (\texttimes) Single CLIP backbone/fixed computation
  & (\texttimes) Single MLLM/fixed computation
  & (\checkmark) LLM cascade/variable computation
  & (\checkmark) Heterogeneous three-stage pipeline \\
\midrule
Sequential routing
  & (\texttimes) No routing/uniform inference
  & (\texttimes) No routing/uniform inference
  & (\texttimes) No routing/uniform inference
  & (\texttimes) No routing/uniform inference
  & (\checkmark) Budget-aware cascade routing
  & (\checkmark) Disagreement-based escalation (primary) \\
\midrule
Cost awareness
  & (\texttimes) Fixed cost per sample
  & (\texttimes) Fixed cost per sample
  & (\texttimes) Fixed cost per sample
  & (\texttimes) Fixed cost per sample
  & (\checkmark) Minimizes LLM API cost
  & (\checkmark) Latency and token-aware cascade \\
\midrule
Multi-model disagreement
  & (\texttimes) Single model/no signal
  & (\texttimes) Single model/no signal
  & (\texttimes) Single model/no signal
  & (\texttimes) Single model/no signal
  & (\texttimes) Not applicable to forgery detection
  & (\checkmark) Artifact/semantic detector disagreement as primary signal \\
\midrule
Generator shift eval.
  & (\texttimes) GAN only, degrades on diffusion
  & (\checkmark) Evaluated across diffusion
  & (\checkmark) Evaluated across multiple models
  & (\checkmark) Modern video generators
  & (\texttimes) Not evaluated on media forgery
  & (\checkmark) Video, GAN, diffusion, SDXL, and FLUX generators \\
\midrule
Media type
  & (\texttimes) Images only
  & (\texttimes) Images only
  & (\texttimes) Images only
  & (\checkmark) Video forgery detection
  & (\texttimes) Text only
  & (\checkmark) Video frames and static images \\
\bottomrule
\end{tabularx}
\end{table*}


\section{Proposed Solution}
This section describes the proposed framework for difficulty-aware model selection. We first formalize the routing problem and define the cost model, then introduce the disagreement-based routing strategy, examine the alternative routing baselines used for comparison, and finally detail how latency is measured across routing strategies.

\subsection{Problem Formulation}

We consider the problem of detecting AI-generated media using a set of heterogeneous detectors with different computational costs and failure modes. Given an input sample $x$, the goal is to select an appropriate subset of detectors to apply such that prediction accuracy is maintained while minimizing computational cost.

Let:
\[
\mathcal{M} = \{f_{\text{art}}, f_{\text{sem}}, f_{\text{vlm}}\}
\]
denote the set of available detectors, corresponding to an artifact-based model, a semantic model, and a high-capacity vision-language model. Each model produces a binary prediction $\hat{y}_i \in \{0,1\}$. A routing policy $\pi(x) \in \{0,1\}$ determines whether the VLM is invoked on input $x$, with the artifact and semantic detectors always applied as
low-cost first-stage filters. We evaluate multiple instantiations of this framework in which the artifact detector is instantiated as a dual-stream FFT probe (primary) or CNNSpot
(baseline), and the semantic detector as CLIP, SigLIP2, or DINOv2.

We report on two complementary cost metrics. Wall-clock latency, measured in milliseconds on a fixed GPU, applies to every detection stage:
\[
C_{\text{lat}}(x, \pi) = \kappa^{\text{lat}}_{\text{art}} + \kappa^{\text{lat}}_{\text{sem}}
+ \mathbbm{1}[\pi(x)=1] \cdot \kappa^{\text{lat}}_{\text{vlm}}(x),
\]
where $\kappa^{\text{lat}}_{\text{vlm}}(x)$ is measured per sample because VLM latency depends on input resolution. The token cost, measured as mean VLM input and output tokens per sample, applies only when the VLM is invoked and reflects the standard unit of cost for commercial VLM APIs~\cite{chen2024frugalgpt}:
\[
C_{\text{tok}}(x, \pi) = \mathbbm{1}[\pi(x)=1] \cdot \bigl(\tau^{\text{in}}_{\text{vlm}}(x)
+ \tau^{\text{out}}_{\text{vlm}}(x)\bigr).
\]
The routing objective is to minimize the expected cost under either metric while maintaining a target accuracy $\alpha$:
\[
\min_{\pi} \; \mathbb{E}_x[C(x, \pi)] \quad \text{s.t.} \quad \text{Accuracy}(\pi) \geq \alpha.
\]
We report both $C_{\text{lat}}$ and $C_{\text{tok}}$ in all evaluated methods, as the deployment scenarios vary between self-hosted (latency-bound) and API-based (token-bound) inference.

\subsection{Disagreement as a Difficulty Signal}

We estimate the input difficulty using disagreement between the artifact-based and semantic detectors:

\begin{equation}
d(x) = \mathbbm{1}\left[\hat{y}_{\text{art}}(x) \neq \hat{y}_{\text{sem}}(x)\right].
\label{eq:disagree}
\end{equation}

When the detectors agree, predictions are highly reliable. Disagreement corresponds to
cases where predictions conflict, indicating that neither signal alone is sufficient and
that additional computation may be required. Unlike confidence-based measures such as
entropy, this signal requires no calibration and shows consistent diagnostic behavior
across generator families. This behavior comes from the complementary failure modes of
both artifact and semantic detectors, whose errors are weakly correlated under
distribution shift. This leads to disagreement reliably identifying the cases where cheap detectors
are not trusted.


\subsection{Disagreement-Based Routing}

Using the disagreement signal, we define a simple routing policy:

\begin{equation}
\pi_{\text{disagree}}(x) =
\begin{cases}
0, & \text{if } d(x) = 0 \\
1, & \text{if } d(x) = 1
\end{cases}
\label{eq:pidisagree}
\end{equation}

where $\pi(x) = 1$ indicates that the VLM is invoked. Inputs where the detectors agree are handled by the cheap stage alone while disagreement cases are escalated to the VLM. This enables selective allocation of computation based on estimated difficulty. Note that outputting the
semantic detector's prediction on all inputs is equivalent to this routing policy without VLM escalation, since the artifact and semantic detectors produce identical predictions on
agreement cases by definition. In-domain, this semantic-only baseline achieves the best
cost-accuracy tradeoff. The cascade's value is under deployment shift, where a stale
semantic probe loses accuracy and VLM escalation on disagreement cases recovers it.

\subsection{Generalized K-Stage Formulation and Three Cascade Cases}
\label{sec:kstage}

The proposed detector is a special case of a $K$-stage cascade. Let
$f_1, \ldots, f_K$ denote detectors with non-decreasing cost
$C_1 \le C_2 \le \cdots \le C_K$ and accuracies $A_1, \ldots, A_K$ on a
fixed evaluation distribution. The useful quantity is the conditional gap
on the inputs that stage $i$ actually routes forward. The global gap
$A_{i+1}-A_i$ is only a coarse summary. Escalation improves accuracy only
when the next stage is better on that routed subset.

In our setting, let $K=3$: $f_1 = f_{\text{art}}$, $f_2 = f_{\text{sem}}$, and
$f_3 = f_{\text{vlm}}$. The routed subset is the disagreement set
$D = \{x : \hat y_{\text{art}}(x) \neq \hat y_{\text{sem}}(x)\}$
(Eq.~\eqref{eq:disagree}). The resulting cascade accuracy is
\begin{equation}
A_{\text{cascade}} = (1-p_D)\, A_{\text{agree}} + p_D\, A_{3\mid D},
\label{eq:decomp}
\end{equation}
where $p_D=P(x\in D)$, $A_{\text{agree}}$ is accuracy on the agreement set,
and $A_{3\mid D}$ is verifier accuracy on $D$. On the agreement set,
$f_1$ and $f_2$ make the same prediction by construction
(Eq.~\eqref{eq:pidisagree}). The gain over the semantic-only detector is
therefore
\begin{equation}
\text{Gain} = A_{\text{cascade}} - A_2
= p_D \left( A_{3\mid D} - A_{2\mid D} \right),
\label{eq:gain}
\end{equation}
where $A_{2\mid D}$ is the accuracy of the semantic detector on the same routed
subset. Thus, a stronger verifier helps only if it is stronger on
disagreement cases. Average accuracy alone does not determine the cascade gain.
This leaves us with three cascade cases.

\textbf{Case 1 (both cheap stages strong, $A_1 \approx A_2 \approx A_3$).}
The artifact and semantic detectors usually agree, so $p_D$ is small. Even
a strong verifier has little effect on end-to-end accuracy because it is
rarely invoked.

\textbf{Case 2 (asymmetric cheap stages, $A_1 \ll A_2$, and $A_2 \gtrsim A_3$).}
The semantic detector is already the better cheap model. Many disagreement
cases are therefore inputs where the artifact detector is wrong and the
semantic detector is right. If $A_{3\mid D} < A_{2\mid D}$, routing replaces
correct semantic predictions with noisier verifier predictions, so
Gain $<0$. In this case, any improvement over the \emph{artifact}-only
baseline comes from the semantic detector, rather than from escalation.

\textbf{Case 3 (both cheap stages degraded, $A_1 \ll A_3$ and $A_2 \ll A_3$).}
Both cheap detectors are unreliable with disagreement being common. If the
verifier remains accurate on $D$, then $A_{3\mid D}$ can substantially
exceed $A_{2\mid D}$. This is the cascade case where Eq.~\eqref{eq:gain} becomes
large and positive, and the gain is caused by routing hard cases to the
verifier.

\textbf{Breakdown condition.} Disagreement is useful only when the two cheap
detectors fail differently and the verifier is better on the routed subset.
If $f_1$ and $f_2$ share a correlated blind spot, they can agree while both
are wrong. Those inputs never enter $D$, so the cascade cannot route them to
the verifier. In the extreme, $p_D \to 0$ and
Eq.~\eqref{eq:decomp} reduces to
$A_{\text{cascade}} \approx A_{\text{agree}}$, even when the agreement set
contains systematic errors. Section~VI maps these cases and this breakdown
condition to the empirical results in Section~V.

\subsection{Baseline Routing Strategies}

To evaluate the effectiveness of disagreement-based routing, we compare two alternative strategies.

\textbf{Entropy-based routing.}  
We use the entropy of the artifact detector as a measure of uncertainty:

\begin{equation}
H(x) = - \sum_{i} p_i(x) \log p_i(x),
\end{equation}

where $p_i(x)$ are the predicted class probabilities. Inputs with entropy above a threshold are escalated to the VLM. This represents a standard confidence-based routing approach.

\textbf{Bandit-based routing.}
We train a contextual bandit model using a logistic regression policy on features derived from detector output, including artifact detector confidence, semantic confidence, and prediction disagreement. The model predicts whether an input should be escalated to the VLM, providing a learned routing baseline.

\subsection{Latency Measurement}

Inference latency for each model component is measured independently using synchronized CUDA execution on a fixed hardware setup. The total latency per sample is the sum of the costs of all models invoked by the routing policy. Artifact and semantic detectors are modeled as additive costs, which is a  conservative estimate since these stages could run in parallel in practice. All latency measurements are intended for relative comparison across routing strategies rather than as absolute deployment guarantees.

\FloatBarrier
\begin{figure*}[h!]
    \centering
    \includegraphics[width=0.7\textwidth]{workflowfinal.png}
    \caption{Comparison of detection workflows. (a) FrugalGPT introduces cost-aware cascade
routing for homogeneous LLM outputs. (b) Our framework adapts this paradigm to
heterogeneous forensic detectors, using artifact-semantic disagreement as the primary
routing signal.}
    \label{fig:workflowdiagram}
\end{figure*}
\FloatBarrier

\section{Experiments}

This section describes the evaluation methodology used to assess the proposed routing framework. We first define the compared methods and evaluation metrics, then summarize implementation details and latency measurements.

\subsection{Evaluation Strategy}
We evaluate our framework against artifact-only and semantic-only detectors applied without routing, an always-escalate pipeline that invokes the verifier on every sample, and entropy- and bandit-based routing baselines. UnivFD~\cite{ojha2023fakedetect} is included
as a published 2023 baseline using CLIP ViT-L/14 features trained on GAN images. We
evaluate two verifiers: Qwen2.5-VL-7B~\cite{bai2025qwen25vltechnicalreport} as the
baseline and GPT-5.5~\cite{openai2026gpt55} as the primary, to test whether a stronger verifier better exploits the disagreement signal.
Evaluation focuses on detection accuracy, balanced accuracy, 
average computational cost, and escalation rate. For video 
datasets, a representative frame is sampled from each clip 
and treated as an image detection task following standard 
evaluation practice~\cite{wen2025busterxmllmpoweredaigeneratedvideo}.



\subsection{Implementation}

The pipeline integrates three heterogeneous detection components. The first stage is an artifact-based detector. The primary artifact detector is a dual-stream frequency-domain probe that combines radial FFT magnitude profiles~\cite{dzanic2020fourier} with DCT high-frequency statistics~\cite{frank2020leveraging, Corvi_2023_ICASSP} and color channel
coherence features, trained as a logistic regression classifier. CNNSpot~\cite{wang2020cnn} is retained as an artifact detector baseline. The second stage is a semantic detector. Visual embeddings are extracted with one of three backbones: CLIP
ViT-B/32~\cite{pmlr-v139-radford21a}, SigLIP2~\cite{tschannen2025siglip2}, or
DINOv2~\cite{oquab2023dinov2}. It's then passed through a logistic regression classifier trained on the first 50\% of each dataset. All three semantic probes use the same training split and evaluation protocol. The third and final stage is an optional VLM verifier which is only invoked 
on disagreement cases.


To utilize GPT-5.5 as the primary image verifier, we prompt it as follows:

\noindent
\begin{minipage}{\linewidth}
\vspace{10pt}
\itshape
``Analyze this image carefully. Determine whether it is a REAL photograph or an AI-generated FAKE image. Consider both possibilities equally. Answer with exactly one word: REAL or FAKE.''
\vspace{10pt}
\end{minipage}

When applied to the real images bundled with AIGIBench (PASCAL VOC 2007--2012
photographs), GPT-5.5 classified 82--97\% of them as \emph{FAKE} despite being
genuine photographs. PASCAL VOC images are archival scans with heavy JPEG
compression that a frontier VLM trained on modern high-quality imagery does not
recognize as real, reducing real-class recall to near zero. We replaced these
images with ImageNet val/nature photographs from the GenImage SD~1.4
partition~\cite{NEURIPS2023_f4d4a021}, which GPT-5.5 correctly identifies as real
and which match the real-image source used to train the SD~1.4 semantic probes.
Any VLM-based evaluation on AIGIBench should substitute the bundled real images
with contemporary photographs to avoid this calibration artifact.


Confidence scores are derived from the logistic probe output for the FFT artifact detector and all semantic detectors. For the CNNSpot baseline, confidence is the maximum softmax probability from the ResNet output. The bandit and entropy baselines use these scores as features. The disagreement rule uses only binary predictions and requires no confidence
calibration.

The routing policy operates on the joint output of the artifact and semantic detectors. 
In the disagreement threshold configuration, escalation is triggered whenever the two first-stage detectors produce conflicting 
predictions. In the entropy threshold configuration, escalation is
triggered when the artifact detector's prediction entropy exceeds $\tau = 0.5$. The 
bandit router applies a logistic regression policy trained on artifact
confidence, semantic confidence, and their disagreement signal to produce 
a binary escalation decision.

We measure deployment latency on an NVIDIA RTX A4500 GPU with synchronized CUDA execution. The FFT probe takes 2.5~ms per sample. Each semantic detector takes 10.4~ms per sample,
giving a cheap-stage latency of 12.9~ms. This represents a 4.4$\times$ reduction in cheap-stage cost compared to the CNNSpot-based configuration (56.2~ms). GPT-5.5 latency on the original benchmarks averages 1{,}595~ms on GenBuster, 1{,}459~ms on SD~1.4, and 1{,}572~ms on BigGAN. On the SDXL and FLUX benchmarks, where images are higher resolution, GPT-5.5 latency averages 3{,}084--3{,}103~ms and 3{,}670--3{,}849~ms respectively. Qwen2.5-VL
latency on the original benchmarks is 190~ms (GenBuster), 83~ms (SD~1.4), and 71~ms
(BigGAN) under always-escalate. We report both latency and token cost throughout our results. Latency comparisons assume sequential FFT and semantic detector execution. Parallel
execution would further reduce cheap-stage cost.

\section{Results}

This section presents empirical results across video, diffusion, and GAN benchmarks. We first describe the datasets and evaluation protocol, then compare detector performance, analyze routing behavior, and finally study the robustness of the proposed signal through ablations.

\subsection{Experimental Setup}

We evaluate on three in-domain benchmarks spanning distinct generator families.
GenBuster-200K-mini~\cite{wen2025busterxmllmpoweredaigeneratedvideo} comprises 2{,}000
samples from four modern AI video generators and 1{,}000 real video clips.
GenImage SD~1.4~\cite{NEURIPS2023_f4d4a021} comprises 6{,}000 Stable Diffusion v1.4
images and 6{,}000 real ImageNet images. GenImage
BigGAN~\cite{NEURIPS2023_f4d4a021} comprises 6{,}000 BigGAN images and 6{,}000 real
ImageNet images. Each sample is represented by a single frame or image following standard
practice~\cite{wen2025busterxmllmpoweredaigeneratedvideo}.

We additionally evaluate under deployment shift on two modern diffusion benchmarks.
Fake images are drawn from AIGIBench~\cite{li2025artificial} SDXL and FLUX1-dev
subsets. Real images are ImageNet val/nature photographs from the GenImage SD~1.4
partition. AIGIBench's bundled PASCAL VOC real images are not used because they
cause systematic VLM calibration failure (Section~IV.B). Each benchmark contains
4{,}000 samples (2{,}000 fake, 2{,}000 real). All semantic probes are trained
exclusively on SD~1.4 data. No SDXL or FLUX images appear in any training split.


For the in-domain benchmarks, each dataset is partitioned into three non-overlapping splits
using stratified sampling. The first split (50\%) trains the semantic probe classifiers,
with a held-out 20\% subset used for probe validation. The remaining 50\% forms the routing
policy training and evaluation pool, from which a held-out 40\% subset is used for all
reported evaluations. No detector or routing policy is evaluated on samples seen during
training. Evaluation sets contain 400 samples for GenBuster and 480 samples each for
SD~1.4 and BigGAN.

We compare FFT-only, CNNSpot-only~\cite{wang2020cnn}, UnivFD~\cite{ojha2023fakedetect},
and semantic-only configurations using CLIP~\cite{pmlr-v139-radford21a},
SigLIP2~\cite{tschannen2025siglip2}, and DINOv2~\cite{oquab2023dinov2}. Routing
baselines include entropy threshold ($\tau = 0.5$), disagreement-based escalation to Qwen2.5-VL-7B or GPT-5.5, and uniform escalation to Qwen2.5-VL. All configurations
are evaluated under identical conditions on the same held-out splits.

All accuracy figures are reported with 95\% bootstrap confidence intervals over 1{,}000
resamples. Balanced accuracy is the mean of fake and real recall. Latency is measured in
milliseconds on an NVIDIA RTX A4500 GPU.

\subsection{Detector Performance}


\textbf{Artifact detectors.} Table~\ref{tab:main_results} compares the dual-stream
FFT probe against CNNSpot across all three in-domain benchmarks. FFT achieves 73.5\%,
74.8\%, and 99.8\% accuracy on GenBuster, SD~1.4, and BigGAN respectively, against
CNNSpot's 54.0\%, 50.4\%, and 72.1\%. This gain comes at 2.5~ms per sample versus
CNNSpot's 45.8~ms, an 18$\times$ reduction in artifact-stage latency. The FFT probe
maintains balanced recall across fake and real classes, while CNNSpot exhibits strong real
bias on out-of-domain generators (10.8\% fake recall on GenBuster).

\textbf{Semantic detectors.} Among the three semantic probes, SigLIP2 achieves the
strongest in-domain accuracy on all benchmarks: 98.5\% on GenBuster, 92.7\% on SD~1.4,
and 99.6\% on BigGAN. CLIP and DINOv2 are competitive with each other but behind
SigLIP2 (92.0\% and 96.3\% respectively on GenBuster, 89.8\% and 89.0\% on SD~1.4, 99.0\%
for both on BigGAN).
All three probes are trained on in-domain data and run at 10.4~ms per sample.

When semantic probes are trained in-domain, the semantic-only baseline achieves the best
cost-accuracy tradeoff on all three benchmarks. Routing configurations that invoke
Qwen2.5-VL on disagreement cases achieve 89--90\% on GenBuster and 84--87\% on SD~1.4,
below the semantic-only baselines (96.3--98.5\% and 89--92.7\% respectively). We expected this since
Qwen provides weak signal on these benchmarks (69.0\%, 60.6\%, 70.4\% under
uniform escalation), and invoking it on the hardest cases cannot recover to the level of a
well-trained in-domain probe. This is the expected signature of \textbf{Case 2} in the K-stage formulation
(Section~III.D). On GenBuster and SD~1.4, the artifact detector is weaker than
the semantic detector, while the verifier is not better than the semantic detector
on the routed disagreement subset. For FFT+SigLIP2 disagreements, SigLIP2's own
accuracy $A_{2\mid D}$ is 95.1\% on GenBuster and 82.6\% on SD~1.4, both well
above Qwen's accuracy on the same subset ($A_{3\mid D}=$ 63.6\% and 60.6\%,
respectively; Table~\ref{tab:verifier_comparison}).

Although the disagreement cascade outperforms the FFT baseline on GenBuster and
SD~1.4 (Table~\ref{tab:main_results}), this gain comes from replacing the
artifact detector with the stronger semantic detector on disagreement cases.
It should not be interpreted as evidence that VLM escalation is helping in this
setting. Escalation becomes useful only in Case 3, where the verifier is stronger
than the cheap detectors on the routed subset.

\textbf{Verifier comparison.} Table~\ref{tab:verifier_comparison} compares Qwen2.5-VL
and GPT-5.5 on the FFT+SigLIP2 disagreement subset, the same subset the primary
cascade escalates throughout this paper, rather than a subset defined by a different
detector pair. On this subset, Qwen2.5-VL is \emph{more} accurate than GPT-5.5 on
GenBuster (63.6\% vs.\ 55.5\%, a $-$8.1 point GPT-5.5 delta). The GenBuster disagreement
cases that FFT and SigLIP2 disagree on turn out to be ones where GPT-5.5 does
comparatively poorly, the opposite of the ordering suggested by a subset built from a
weaker artifact detector. On SD~1.4, GPT-5.5's advantage over Qwen widens substantially
on the correct subset (77.3\% vs.\ 60.6\%, $+$16.7 points, versus a $+$5.0-point gap on
the old subset). Because FFT and SigLIP2 agree on nearly all BigGAN inputs, the disagreement subset
contains only three samples, offering little basis for conclusions about relative
verifier accuracy. This dataset-dependent reversal, in which GPT-5.5 is not uniformly stronger than Qwen
even when restricted to the disagreement cases the cascade actually escalates, is itself
evidence that verifier choice should be validated per-deployment rather than assumed
from a single benchmark, and reinforces that Case~2/Case~3 identity depends on which
verifier is paired with which detector pair on which dataset.

\textbf{UnivFD baseline.} UnivFD~\cite{ojha2023fakedetect}, trained on GAN images using
CLIP ViT-L/14, achieves zero fake recall on GenBuster and near-zero on SD~1.4 (5.9\%),
matching chance accuracy. On BigGAN, where the test distribution matches its training
domain, it achieves 91.5\%. GAN-trained detectors fail categorically on video generators
and diffusion images regardless of backbone strength. The advantage of trainable semantic
probes over fixed pretrained ones is consistent across all evaluated backbones.



\begin{figure*}[t]
    \centering
    \includegraphics[width=.7\textwidth]{plot_accuracy_comparison.png}
   \caption{Detector accuracy across the three in-domain benchmarks. Results compare artifact detectors, semantic probes, and the FFT+SigLIP2+Qwen disagreement cascade.}

    \label{fig:accuracy_comparison}
\end{figure*}


\begin{table}[h]
\centering
\footnotesize
\setlength{\tabcolsep}{2.5pt}
\caption{Detection accuracy and latency across the three in-domain benchmarks. Confidence intervals reported later in supplementary material.}
\label{tab:main_results}
\begin{tabular}{lcccccc}
\toprule
 & \multicolumn{2}{c}{GenBuster} & \multicolumn{2}{c}{SD~1.4} & \multicolumn{2}{c}{BigGAN} \\
\cmidrule(lr){2-3} \cmidrule(lr){4-5} \cmidrule(lr){6-7}
Method & Acc. & Lat. & Acc. & Lat. & Acc. & Lat. \\
\midrule
FFT          & 0.735 &  2.5 & 0.748 &  2.5 & 0.998 &  2.5 \\
CNNSpot      & 0.540 & 45.8 & 0.504 & 45.8 & 0.721 & 45.8 \\
UnivFD       & 0.512 & 56.2 & 0.529 & 56.2 & 0.915 & 56.2 \\
CLIP         & 0.920 & 10.4 & 0.898 & 10.4 & 0.990 & 10.4 \\
SigLIP2      & 0.985 & 10.4 & 0.927 & 10.4 & 0.996 & 10.4 \\
DINOv2       & 0.963 & 10.4 & 0.890 & 10.4 & 0.990 & 10.4 \\
Entropy      & 0.527 & 17.9 & 0.525 & 14.3 & 0.733 & 12.9 \\
Casc. (C)    & 0.900 & 55.1 & 0.865 & 36.9 & 0.998 & 13.3 \\
Casc. (S)    & 0.898 & 57.9 & 0.867 & 31.9 & 0.998 & 13.2 \\
Casc. (D)    & 0.890 & 58.4 & 0.842 & 32.5 & 0.998 & 13.5 \\
Always esc.  & 0.690 & 190.2 & 0.606 & 83.2 & 0.704 & 70.6 \\
\bottomrule
\end{tabular}
\end{table}
% --- Table: Verifier comparison (E4) ---
\begin{table}[h]
\centering
\small
\setlength{\tabcolsep}{3pt}  % reduced from 4pt to 3pt
\caption{Verifier accuracy on the FFT+SigLIP2 disagreement subset. Improvement is reported relative to Qwen2.5-VL-7B.}
\label{tab:verifier_comparison}
\begin{tabular}{llcccc}
\toprule
Dataset & Verifier & Acc. & Improvement & $n$ & Lat. (ms) \\
\midrule
GenBuster & Qwen2.5-VL-7B  & 0.636 & ---        & 110 & 174 \\
GenBuster & GPT-5.5        & 0.555 & $-$0.081   & 110 & 2443 \\
\midrule
SD~1.4    & Qwen2.5-VL-7B  & 0.606 & ---        & 132 & 69 \\
SD~1.4    & GPT-5.5        & 0.773 & $+$0.167   & 132 & 2492 \\
\midrule
BigGAN    & Qwen2.5-VL-7B  & 0.667 & ---        &   3 & 53 \\
BigGAN    & GPT-5.5        & 0.667 & $+$0.000   &   3 & 2023 \\
\bottomrule
\end{tabular}
\end{table}

% --- Table: E8 SDXL/FLUX deployment shift main results ---
\begin{table}[h]
\centering
\footnotesize
\setlength{\tabcolsep}{3pt}
\caption{Deployment-shift results on SDXL and FLUX. Semantic probes and FFT are trained on SD~1.4 and evaluated without retraining. Gain denotes cascade accuracy minus stale semantic accuracy.}
\label{tab:e8_results}
\begin{tabular}{lcccc}
\toprule
 & \multicolumn{2}{c}{SDXL} & \multicolumn{2}{c}{FLUX} \\
\cmidrule(lr){2-3} \cmidrule(lr){4-5}
Sem. & Stale FFT & Stale Sem. & Stale FFT & Stale Sem. \\
\midrule
CLIP    & 0.622 & 0.872 & 0.578 & 0.920 \\
SigLIP2 & 0.622 & 0.847 & 0.578 & 0.909 \\
DINOv2  & 0.622 & 0.855 & 0.578 & 0.815 \\
\midrule
 & \multicolumn{2}{c}{Qwen (Acc / Gain)} & \multicolumn{2}{c}{GPT-5.5 (Acc / Gain)} \\
\cmidrule(lr){2-3} \cmidrule(lr){4-5}
CLIP    & 0.826 / $-$0.047 & 0.893 / $+$0.021 & 0.779 / $-$0.140 & 0.914 / $-$0.006 \\
SigLIP2 & 0.817 / $-$0.030 & 0.881 / $+$0.034 & 0.779 / $-$0.130 & 0.913 / $+$0.004 \\
DINOv2  & 0.824 / $-$0.031 & 0.891 / $+$0.036 & 0.734 / $-$0.081 & 0.864 / $+$0.049 \\
\bottomrule
\end{tabular}
\end{table}


\subsection{Routing Analysis}

In this section we analyze whether artifact-semantic disagreement between the FFT artifact detector and
the semantic probe provides a reliable routing signal, and how selective escalation
affects the cost-accuracy tradeoff across in-domain and deployment-shift conditions.

Figure~\ref{fig:agree_disagree} shows disagreement rates across all five evaluation
datasets. On in-domain benchmarks, the rate ranges from 26--28\%. On BigGAN it collapses
to under 1\%, consistent with the FFT probe near-saturating at 99.8\% accuracy and
leaving almost no room for artifact-semantic conflict. Under deployment shift on SDXL and
FLUX, the rate rises to 35--42\%, marking the conditions where the stale semantic probe
loses accuracy and VLM escalation provides gain.


\begin{figure*}[t]
    \centering
    \includegraphics[width=.7\textwidth]{plot_fig_c_disagreement_rate.png}
    \caption{Artifact-semantic disagreement rates across the five evaluation datasets. Disagreement defines the subset routed to the verifier.}
    \label{fig:agree_disagree}
\end{figure*}

When the FFT artifact detector and semantic probe agree, the agreement subset achieves
above 96\% accuracy across all three in-domain benchmarks (Table~\ref{tab:cheap_path}).
When they disagree, FFT accuracy on the disagreement subset drops to under 18\% on
GenBuster and SD~1.4, and neither detector is reliable on its own. This asymmetry
confirms that artifact-semantic disagreement is a stable routing signal without requiring
confidence calibration.

Table~\ref{tab:escalation} breaks down escalation behavior by case type on GenBuster for
the FFT+SigLIP2 configuration. The disagreement rule escalates all 110 disagreement cases
and bypasses all 290 agreement cases by construction. FFT is correct on 289 of 290
agreement-subset samples. The single case where both detectors agree but FFT is wrong
represents the fundamental limit of the routing signal. When both cheap detectors fail
simultaneously, no artifact-semantic disagreement can flag the error. At 99.7\% agreement-
subset accuracy, this limit affects under 0.3\% of all inputs.

\begin{table}[h]
\centering
\caption{Escalation behavior by case type on GenBuster-200K-mini (FFT+SigLIP2). Disagreement captures 105 of the 106 total FFT failures. The sole miss occurs when both detectors make the same incorrect prediction, which disagreement-based routing cannot detect by construction.}
\label{tab:escalation}
\begin{tabular}{lcc}
\toprule
Case & $n$ & Disagree esc. \\
\midrule
FFT correct, models agree    & 289 & 0.00 \\
FFT correct, models disagree &   5 & 1.00 \\
FFT wrong, models agree      &   1 & 0.00 \\
FFT wrong, models disagree   & 105 & 1.00 \\
\bottomrule
\end{tabular}
\end{table}

Table~\ref{tab:per_generator} breaks down accuracy by generator on GenBuster.
CNNSpot achieves near-zero to modest fake recall across the four fake-video generators
(0--17.8\%), confirming that GAN-trained detectors generalize poorly to modern video
content. Disagreement-routed escalation recovers accuracy across generators but falls
below the semantic-only baseline, consistent with Qwen's limited accuracy on video
frames.


\begin{table}[h]
\centering
\caption{Per-generator detection accuracy on GenBuster-200K-mini for artifact, semantic, and disagreement-routed cascade models.}
\label{tab:per_generator}
\begin{tabular}{lrcccc}
\toprule
Generator & $n$ & CNNSpot & Cascade & CLIP \\
\midrule
CogVideoX     & 73  & 0.178 & 0.822 & 0.849 \\
EasyAnimate   & 73  & 0.041 & 0.918 & 0.904 \\
HunyuanVideo  & 16  & 0.000 & 0.875 & 0.813 \\
LTX-Video     & 33  & 0.152 & 0.879 & 0.788 \\
\midrule
Real   & 205 & 0.951 & 0.922 & 0.980 \\
\bottomrule
\end{tabular}
\end{table}


The retrained bandit router, trained on the same FFT+SigLIP2+Qwen features as the disagreement threshold configuration, converges to performance that matches the disagreement rule on all three datasets: SD~1.4 (0.867 vs.\ 0.867) and BigGAN (0.998 vs.\ 0.998) exactly, and GenBuster (0.917 vs.\ 0.898) closely. Section~VI.C interprets this convergence.



Relative to always-escalate to Qwen (190.2~ms per sample, Table~\ref{tab:main_results}),
the disagreement cascade (FFT+SigLIP2) reduces mean per-sample latency by 70\% on
GenBuster (57.9~ms), and the semantic-only baseline alone reduces it by 95\% (10.4~ms). The FFT-based disagreement rule consistently improves over the FFT baseline on
GenBuster and SD~1.4 while matching it on BigGAN
(Table~\ref{tab:main_results}).


\begin{figure*}[h]
    \centering
    \includegraphics[width=.85\textwidth]{plot_cost_accuracy_all.png}
    \caption{Cost-accuracy tradeoff for each semantic backbone under in-domain shown left and deployment-shift evaluation shown on right.}
    \label{fig:cost_accuracy}
\end{figure*}


\subsection{Ablation Studies}
The ablation studies characterize cheap-path coverage, robustness of the
disagreement-based routing signal, and cascade behavior under cross-family probe
transfer.


\subsubsection{Cheap-Path Coverage}

Table~\ref{tab:cheap_path} quantifies cheap-path coverage for the FFT+SigLIP2
configuration. On every dataset, at least 74\% of samples fall in the agreement subset where FFT accuracy exceeds 96\%, and escalation is limited to the remaining disagreement
cases. On BigGAN, the agreement rate reaches 99.4\% and cheap-path accuracy is 100\%,
leaving only 3 samples for VLM escalation.

\begin{table}[h]
\centering
\caption{Agreement-subset coverage and accuracy for the FFT+SigLIP2 configuration. Remaining samples form the disagreement subset routed to the verifier.}

\label{tab:cheap_path}
\begin{tabular}{lcccc}
\toprule
Dataset & Agree & Cheap Acc & Disagree & Esc. \\
        & $n$   &           & $n$      & rate \\
\midrule
GenBuster & 290 & 0.997 & 110 & 0.275 \\
SD~1.4    & 348 & 0.966 & 132 & 0.275 \\
BigGAN    & 477 & 1.000 &   3 & 0.006 \\
\bottomrule
\end{tabular}
\end{table}

\subsubsection{Cheap-Trigger Ablation}
\label{sec:cheap_trigger}

Since the FFT detector contributes no independent prediction on agreement cases by
construction (Section~VI.D), we test whether its value comes from the specific
disagreement signal or from escalating a matched-rate subset of inputs. We compare the
disagreement-routed cascade with a random-trigger baseline. For each dataset, the random
baseline escalates the same fraction of inputs as the FFT+SigLIP2 disagreement rule, but
chooses those inputs uniformly at random. We repeat this process over 2{,}000 draws and
report a 95\% confidence interval for the random baseline's cascade accuracy.

Table~\ref{tab:cheap_trigger} shows that disagreement routing significantly outperforms
random triggering on SD~1.4 and BigGAN. In these datasets, the disagreement subset
identifies inputs where the semantic probe is less reliable than it is on average, so
routing those specific cases to Qwen improves accuracy more than escalating a random
subset of the same size. On GenBuster, the two methods are statistically
indistinguishable. The semantic probe remains accurate even on disagreement cases
(95.5\% versus 98.5\% overall), so disagreement does not isolate a substantially harder
subset. This suggests that FFT provides a useful routing signal on SD~1.4 and BigGAN,
while its added value is not established on GenBuster, where escalation is already least
useful (Case 2, Section~V.B).

\begin{table}[h]
\centering
\caption{Disagreement-routed cascade vs.\ a matched-rate random-trigger baseline (2,000 draws, 95\% CI on the random-trigger mean).}
\label{tab:cheap_trigger}
\footnotesize
\setlength{\tabcolsep}{3pt}
\begin{tabular}{l c c c c}
\toprule
Dataset & Esc. Rate & Disag. Casc. & Random Casc. [95\% CI] & Beats Random? \\
\midrule
GenBuster & 0.275 & 0.898 & 0.904 [0.880, 0.925] & No \\
SD~1.4 & 0.275 & 0.867 & 0.839 [0.817, 0.860] & Yes \\
BigGAN & 0.006 & 0.998 & 0.994 [0.990, 0.996] & Yes \\
\bottomrule
\end{tabular}
\end{table}

\subsubsection{Probe-Degradation Sweep}
\label{sec:limited_data}

We next test how cascade gain changes as the semantic probe becomes stronger. We hold
the FFT artifact detector fixed and vary the fraction of labeled training data available
to the semantic probe (1\%, 5\%, 10\%, 25\%, 50\%, 100\%). At each fraction, we compare
semantic-only accuracy with FFT+semantic disagreement routing using Qwen as the verifier.
Table~\ref{tab:limited_data} reports the results of SD~1.4 with SigLIP2 and DINOv2
backbones.

\begin{table}[h]
\centering
\caption{Probe-degradation sweep on SD~1.4. Cascade gain is reported relative to the semantic-only probe at the same training fraction.}
\label{tab:limited_data}
\small
\begin{tabular}{lcccc}
\toprule
Backbone & Train frac. & Sem.-only & Cascade & Gain \\
\midrule
SigLIP2  & 1\%   & 0.740 & 0.794 & $+$0.054 \\
SigLIP2  & 5\%   & 0.842 & 0.833 & $-$0.008 \\
SigLIP2  & 10\%  & 0.898 & 0.854 & $-$0.044 \\
SigLIP2  & 25\%  & 0.908 & 0.863 & $-$0.046 \\
SigLIP2  & 50\%  & 0.919 & 0.869 & $-$0.050 \\
SigLIP2  & 100\% & 0.927 & 0.867 & $-$0.060 \\
\midrule
DINOv2   & 1\%   & 0.669 & 0.769 & $+$0.100 \\
DINOv2   & 5\%   & 0.781 & 0.798 & $+$0.017 \\
DINOv2   & 10\%  & 0.808 & 0.810 & $+$0.002 \\
DINOv2   & 25\%  & 0.829 & 0.808 & $-$0.021 \\
DINOv2   & 50\%  & 0.852 & 0.821 & $-$0.031 \\
DINOv2   & 100\% & 0.890 & 0.842 & $-$0.048 \\
\bottomrule
\end{tabular}
\end{table}

Both backbones move from positive to negative cascade gain as the semantic probe improves.
For SigLIP2, gain is positive at 1\% training data and becomes negative by 5\%. For
DINOv2, gain is positive through 10\% and becomes negative by 25\%. This matches
Eq.~\eqref{eq:gain}: escalation helps when Qwen is stronger than the semantic probe on
the disagreement subset, but hurts once the semantic probe becomes the better model on
those routed cases.

GenBuster does not show the same crossover within the tested range. Even with only
1\% of GenBuster's training data (16 labeled examples), the semantic probe already
outperforms the FFT detector, so cascade gain remains negative. BigGAN also does not
cross over, but for a different reason: FFT already exceeds 99\% accuracy and the
disagreement rate is near zero, keeping the setting in Case~1. The sweep therefore
shows the expected movement between Case~2 and Case~3 as semantic-probe quality changes,
while also showing that the transition depends on the dataset and routed subset.

\subsubsection{Cross-Dataset Bandit Transfer}

Table~\ref{tab:transfer} reports cross-dataset bandit transfer accuracy using the same
FFT+SigLIP2+Qwen routing features as the disagreement rule. The in-domain diagonal
entries closely match the disagreement rule: 0.867 on SD~1.4, 0.998 on BigGAN, and
0.917 versus 0.898 on GenBuster. Off-diagonal transfer ranges from 0.758 to 0.998.
BigGAN is the easiest transfer target, while SD~1.4 is the weakest target, consistent
with Qwen's noisier signal on that benchmark.

The disagreement rule achieves 0.898, 0.998, and 0.867 on GenBuster, BigGAN, and
SD~1.4 respectively, matching or exceeding every off-diagonal bandit transfer result.
This is notable because the disagreement rule has no trainable parameters and requires
no per-dataset retraining. The in-domain bandit reproduces the disagreement rule exactly
on SD~1.4 and BigGAN, and performs only slightly better on GenBuster, suggesting that
the learned policy mainly recovers the same routing signal rather than discovering a
stronger one.



\begin{table}[h]
\centering
\caption{Cross-dataset contextual bandit transfer accuracy using FFT+SigLIP2+Qwen routing features.}
\label{tab:transfer}
\begin{tabular}{lccc}
\toprule
Train $\rightarrow$ Test & GenBuster & BigGAN & SD14 \\
\midrule
GenBuster & \textbf{0.917} & 0.998 & 0.867 \\
BigGAN    & 0.843 & \textbf{0.998} & 0.758 \\
SD14      & 0.897 & 0.998 & \textbf{0.867} \\
\midrule
Disagreement rule (no training) & 0.898 & 0.998 & 0.867 \\
\bottomrule
\end{tabular}
\end{table}


\subsubsection{Sensitivity to Semantic Probe Quality}

The disagreement signal assumes that both cheap detectors provide useful predictions.
When one detector falls to near-random accuracy, disagreement no longer identifies hard
cases reliably. Cross-family experiments show that semantic probes trained on one
generator family transfer poorly to distant families. For example, a GenBuster-trained
probe evaluated on SD~1.4 reaches only 0.597 accuracy, close to chance. Across
off-diagonal source-target pairs, semantic-probe accuracy ranges from 0.505 to 0.598,
showing that cross-family probe transfer is unreliable regardless of source domain.

This limits how the cascade should be deployed. Applying it to a new generator family
requires an in-domain semantic probe, or at least a probe that remains meaningfully above
chance on the target distribution. The disagreement rule adds no trainable parameters,
but it still depends on both cheap detectors that carry a usable signal.


\subsubsection{Cross-Family Deployment}
\label{sec:cross_family_deployment}

Cross-family deployment tests whether the cascade helps when the semantic probe is
applied outside its training family. A SigLIP2 probe trained on dataset A is evaluated
on dataset B without retraining. We compare the stale probe alone, an FFT-only detector
trained in-domain on the target dataset, and the disagreement cascade with Qwen as the
verifier (Table~\ref{tab:cross_family_deployment}). For reference, in-domain SigLIP2
accuracy is 0.985 on GenBuster, 0.927 on SD~1.4, and 0.996 on BigGAN, showing the scale
of the cross-family degradation.

The cascade improves over the stale SigLIP2 probe in all six pairings, with a mean gain
of 17.5 points. This gain is misleading if the stale probe is the only baseline. Against
FFT-only, the cascade improves in only one of six pairings, with a mean change of
$-$10.2 points. The lone positive result, BigGAN$\to$GenBuster ($+$0.003), is not
statistically distinguishable from zero under a paired bootstrap (95\% CI
$[-0.037,+0.043]$). The SD~1.4$\to$GenBuster difference is also indistinguishable from
zero ($-$0.005, 95\% CI $[-0.048,+0.040]$). The remaining four comparisons are
significantly negative.

These results illustrate \textbf{Case 2} from Section~III.D. The transferred semantic
probe is weak, so the cascade appears to recover accuracy relative to that stale probe.
However, many disagreement cases are inputs that the in-domain FFT detector already
classifies correctly. Escalating those cases to Qwen replaces correct FFT predictions
with a less reliable verifier prediction. Under this cross-family setting, a strong
in-domain artifact detector is therefore preferable to escalating disagreements from a
poorly transferred semantic probe.

This setting differs from the SD~1.4-to-SDXL/FLUX deployment-shift experiment in
Table~\ref{tab:e8_results}. Here, the semantic probe is transferred across generator
families and Qwen is used as the verifier. In Table~\ref{tab:e8_results}, the semantic
probe is trained on SD~1.4 and evaluated on newer diffusion models from the same lineage,
with GPT-5.5 as verifier. The positive cascade result in that setting is therefore a
different Case~3 configuration.

\begin{table}[h]
\centering
\caption{Cross-family deployment results. Cascade uses a cross-family SigLIP2 probe with Qwen escalation on disagreement, while FFT is trained in-domain on each evaluation dataset. $\Delta_{\text{stale}}$ denotes improvement over the stale SigLIP2 probe, and $\Delta_{\text{FFT}}$ denotes improvement over FFT-only. $^\dagger$ marks $\Delta_{\text{FFT}}$ cells whose paired-bootstrap 95\% CI includes zero (10,000 resamples per cell).}
\label{tab:cross_family_deployment}
\footnotesize
\setlength{\tabcolsep}{3pt}
\begin{tabular}{llccccc}
\toprule
Train & Test & FFT & Stale SigLIP2 & Casc. & $\Delta_{\text{stale}}$ & $\Delta_{\text{FFT}}$ \\
\midrule
GenBuster & SD~1.4 & 0.748 & 0.552 & 0.654 & +0.102 & -0.094 \\
GenBuster & BigGAN & 0.998 & 0.554 & 0.787 & +0.233 & -0.210 \\
SD~1.4 & GenBuster & 0.735 & 0.547 & 0.730 & +0.182 & -0.005$^\dagger$ \\
SD~1.4 & BigGAN & 0.998 & 0.613 & 0.833 & +0.221 & -0.165 \\
BigGAN & GenBuster & 0.735 & 0.517 & 0.738 & +0.220 & +0.003$^\dagger$ \\
BigGAN & SD~1.4 & 0.748 & 0.515 & 0.604 & +0.090 & -0.144 \\
\midrule
\multicolumn{2}{l}{Mean} & & & & +0.175 & -0.102 \\
\bottomrule
\end{tabular}
\end{table}

\begin{figure*}[t]
    \centering
    \includegraphics[width=0.7\textwidth]{plot_fig_a_semantic_degradation.png}
    \caption{Semantic probe accuracy under in-domain (solid) and deployment-shift evaluation (striped) on SDXL and FLUX.}
    \label{fig:semantic_degradation}
\end{figure*}

\subsubsection{Distant Generator Shift: BigGAN}
\label{sec:biggan_distant_shift}

SDXL and FLUX are diffusion models from the same broad lineage as SD~1.4, the dataset
used to train the stale FFT and semantic probes. To test a more distant shift, we apply
the SD~1.4-trained FFT and SigLIP2 probes zero-shot to GenImage BigGAN, a GAN-based
generator. This differs from the cross-family deployment experiment
(Section~V.D.6), where FFT is trained in-domain on the target dataset. Here, both cheap
detectors are stale, matching the SDXL/FLUX deployment-shift setup.

The shift to BigGAN causes much stronger degradation than SDXL or FLUX. Stale FFT
accuracy drops to 38.3\%, and stale SigLIP2 drops to 62.1\%, below the corresponding
SDXL/FLUX ranges of 57.8--62.2\% and 81.5--92.0\%
(Figure~\ref{fig:semantic_degradation}). The disagreement rate is 29.3\%. On the
disagreement subset, stale FFT falls to 9.2\% accuracy, while stale SigLIP2 remains
comparatively strong at 90.8\%. This produces an asymmetric Case~2 pattern rather than
the joint cheap-detector degradation needed for Case~3.

We evaluate GPT-5.5, Qwen2.5-VL-7B, and Qwen2.5-VL-72B (4-bit quantized) on a
stratified random sample of 1{,}500 disagreement cases. All three verifiers produce a 
negative cascade gain relative to the stale semantic probe (Table~\ref{tab:biggan_shift}).
The reason follows directly from Eq.~\eqref{eq:gain}. None of the verifiers exceeds
stale SigLIP2's 90.8\% accuracy on the routed subset, so escalation replaces a
comparatively reliable semantic prediction with a weaker verifier prediction. This
matches the Case~2 failure mode from the cross-family deployment experiment
(Section~V.D.6), rather than the Case~3 recovery observed on SDXL/FLUX.

We treat this as a substantive negative result. The positive SDXL/FLUX finding does
not extend to a generator shift of this size. A larger distribution shift can still
hurt escalation if the verifier is weaker than the best cheap detector on the
disagreement subset. Increasing model scale does not remove this issue:
Qwen2.5-VL-72B (4-bit quantized) reaches only 46.7\% accuracy on the BigGAN
disagreement subset, below Qwen2.5-VL-7B's 68.9\%, due to a strong real-class bias.

\begin{table}[h]
\centering
\caption{Cascade gain over the stale semantic probe on the BigGAN distant-shift disagreement subset ($n=1{,}500$), paired-bootstrap 95\% CI over 10,000 resamples. All three verifiers give significantly negative gain.}
\label{tab:biggan_shift}
\footnotesize
\setlength{\tabcolsep}{3pt}
\begin{tabular}{l c c c}
\toprule
Verifier & Acc. on Disagree & Cascade Acc. & Gain [95\% CI] \\
\midrule
GPT-5.5 & 0.801 & 0.590 & -0.031 [-0.041, -0.022] \\
Qwen2.5-VL-7B & 0.689 & 0.557 & -0.064 [-0.074, -0.054] \\
Qwen2.5-VL-72B & 0.467 & 0.492 & -0.129 [-0.140, -0.119] \\
\bottomrule
\end{tabular}
\end{table}


\section{Discussion}

This section interprets the empirical results presented in Section V. We first examine
why artifact-semantic disagreement separates easy from hard inputs, then map the
three K-stage cascade cases from Section~III.D onto the empirical results and assess whether
learned routing improves on the disagreement signal, then characterize the role of the
semantic probe and VLM verifier under in-domain and deployment-shift conditions, and
finally discuss generator distribution shift, limitations, and practical deployment
guidelines.



\subsection{Artifact-Semantic Disagreement as a Routing Signal}

The FFT artifact detector achieves 99.7\%, 96.6\%, and 100\% accuracy on agreement
cases across the three benchmarks and drops to 4.5\%, 17.4\%, and 66.7\% on
disagreement cases. Agreement between the artifact and semantic detectors is a
near-perfect predictor of an easy sample, and disagreement predicts difficulty with
equal reliability. This pattern holds across fundamentally different generator
architectures: video-based diffusion generators (GenBuster), image diffusion models
(SD~1.4), and GAN-based generators (BigGAN).

The complementarity arises from the signal types each detector uses. The FFT probe
captures spectral regularities introduced during synthesis, including upsampling grids
and DCT coefficient patterns. Semantic probes encode high-level scene content learned
from natural image-text distributions. These signal types are largely orthogonal.
When both detectors agree, at least one type of signal is unambiguous and the sample
is easy to classify. When they disagree, neither is sufficient alone, indicating a
sample that requires higher-level reasoning. The disagreement rule routes only these
ambiguous cases to the VLM verifier without requiring a learned policy or calibrated
confidence scores.

BigGAN produces fewer than 1\% disagreement cases in-domain because GAN upsampling
artifacts are unambiguous in the frequency domain. The FFT probe resolves nearly all
BigGAN samples at the cheap stage. On GenBuster and SD~1.4, modern diffusion and
video generators produce 26--28\% disagreement rates, marking the case where VLM
escalation recovers accuracy that neither cheap detector can provide alone.

\subsection{The Three Cascade Cases in Practice}
\label{sec:three_cases}

The K-stage formulation (Section~III.D) predicts three qualitative cascade cases, and
the empirical results exhibit all three. \textbf{Case 1} occurs on BigGAN in-domain.
FFT alone reaches 99.8\% accuracy and disagreement falls below 1\%, so the cascade adds
latency with almost no change in accuracy (Table~\ref{tab:main_results}).

\textbf{Case 2} occurs when the cheap detectors are asymmetric. This describes the
in-domain GenBuster and SD~1.4 results (Section~V.B), and appears more clearly in the
cross-family deployment experiment (Table~\ref{tab:cross_family_deployment},
Section~V.D.6). A poorly transferred SigLIP2 probe creates a large apparent gain over
the stale probe, but the cascade beats FFT-only in only one of six pairings. That one
positive difference is statistically indistinguishable from zero. The BigGAN distant
shift (Section~V.D.7) is another Case~2 instance: stale FFT and stale SigLIP2 degrade
asymmetrically, stale SigLIP2 remains strong on the disagreement subset (90.8\%), and
all tested verifiers, including GPT-5.5, underperform it there.

\textbf{Case 3} occurs on the SD~1.4-to-SDXL/FLUX shift with GPT-5.5 as verifier
(Table~\ref{tab:e8_results}). Both cheap detectors degrade under this shift, and
GPT-5.5 is strong enough on the routed disagreement subset to recover up to 4.9
percentage points over the stale semantic probe, with positive gain in five of six
configurations. The exception is FLUX with CLIP, where the cascade is slightly worse
than the stale semantic probe. This case is verifier-dependent. Replacing GPT-5.5 with
Qwen2.5-VL-7B on the same SDXL/FLUX disagreement subsets makes all six cascade gains
negative ($-$3.0 to $-$14.0 points), turning the same detector setting into a
Case~2-like outcome.

Model scale alone does not determine the case. On SDXL with SigLIP2, Qwen2.5-VL-72B
(4-bit quantized) nearly matches GPT-5.5 on the disagreement subset (89.8\% vs.\
90.0\% accuracy, with cascade gains of $+$0.033 and $+$0.034). On FLUX with SigLIP2,
however, the same 72B verifier behaves like Qwen2.5-VL-7B rather than GPT-5.5
(58.5\% disagreement-subset accuracy versus GPT-5.5's 90.9\%, with cascade gain
$-$0.130). This shows that verifier strength depends on calibration to the target
generator's artifacts, not on model scale alone. Both cheap detectors being weak is
necessary but not sufficient for Case~3. The verifier must also be stronger than the
stronger cheap detector on the routed subset, as Eq.~\eqref{eq:gain} requires.

The probe-degradation sweep (Table~\ref{tab:limited_data}) demonstrates the
Case~2-to-Case~3 transition within a single controlled experiment. On SD~1.4 with
DINOv2, cascade gain is $-$4.8 points at 100\% training data and increases to $+$10.0
points at 1\% training data. SigLIP2 shows the same sign change between 1\% and 5\%
training data. On GenBuster, gain remains negative at every tested fraction because even
a 1\%-trained semantic probe already exceeds FFT accuracy. This shows that the transition
depends on the relative strength of the cheap detectors on the routed subset, not on the
training-data quantity alone.

The \textbf{breakdown condition} from Section~III.D occurs when the two cheap detectors
share a correlated blind spot. Table~\ref{tab:escalation} shows this directly on
GenBuster. The single case where FFT is wrong and the semantic probe agrees with it is
invisible to disagreement-based routing by construction. This represents the fundamental
limit of the routing rule, which affects less than 0.3\% of inputs in this experiment.

\subsection{Limitations of Learned Routing}

The bandit router uses the same FFT+SigLIP2+Qwen features as the disagreement rule
(Table~\ref{tab:transfer}). Its performance closely matches the parameter-free rule
across all three datasets. On SD~1.4, the bandit achieves 0.867 accuracy, identical to
the disagreement rule. On BigGAN, both achieve 0.998. On GenBuster, the bandit reaches
0.917, close to the disagreement rule's 0.898.

A learned router with 11 detector-derived features and an explicit latency reward does
not provide a consistent improvement over binary disagreement. This suggests that the
dominant routing signal is already captured by whether the artifact and semantic
detectors agree, and that the learned policy mainly recovers the same decision boundary
rather than discovering a stronger routing rule.


\subsection{Role of the Semantic Probe}

The semantic probe (SigLIP2) achieves 98.5\%, 92.7\%, and 99.6\% accuracy at
12.9~ms cheap-stage latency across the three in-domain benchmarks. This is
Pareto-optimal in-domain where not invoking the VLM improved accuracy over the
semantic probe alone.

The role of the FFT detector in the cascade is diagnostic, not predictive. In agreement
cases, the FFT and semantic probe produce identical predictions by construction. Any
routing policy that escalates only disagreement cases to the VLM reduces to outputting
the semantic probe's prediction on agreement cases unchanged. The FFT detector
 does not contribute to independent prediction on these samples. It serves as a cheap signal
whose agreement with the semantic probe identifies easy inputs, routing only the
residual disagreement cases to the VLM.

This clarifies when the cascade adds value. In-domain, the semantic probe is accurate
enough that VLM escalation does not improve accuracy and only adds latency. Under
deployment shift, when a stale probe loses accuracy on unseen generators, disagreement
rises and targeted VLM escalation recovers it. The SDXL and FLUX experiments confirm
this (Table~\ref{tab:e8_results}). Semantic probes do not transfer across generator families, consistent with our cross-family deployment experiments, so production deployment on a new generator
family requires in-domain probe training. The disagreement rule itself needs no retraining.

\subsection{VLM Domain Sensitivity}

Using the prompt in Section~IV.B with REAL/FAKE logit extraction, Qwen2.5-VL achieves
balanced accuracy of 50.8\% on GenBuster, 62.3\% on SD~1.4, and 73.4\% on BigGAN when
applied uniformly to all inputs. This improves substantially over text-decoded Qwen
configurations evaluated during development, which produced near-zero fake recall on
static image datasets. The difference comes from using the REAL/FAKE logit gap at
generation time instead of parsing decoded output tokens. With standard greedy decoding,
the model's text outputs were biased toward ``REAL'' even when its logits separated the
two classes. This is a methodological point for VLM-based binary classification, not only
for this detector.

Even with logit extraction, Qwen2.5-VL remains weaker than the trained semantic probes on
all three in-domain benchmarks. Selective escalation can still make Qwen useful relative
to uniform invocation by applying it only to the artifact-semantic disagreement subset.
On GenBuster, for example, the disagreement cascade reaches 89.8\% accuracy, while
uniform Qwen reaches only 69.0\%. This does not mean Qwen improves over the semantic-only
baseline. Instead, it shows that a weak verifier can still add value relative to weaker
cheap-stage baselines when routed to a targeted subset.

Qwen's accuracy on routed disagreement cases varies substantially across settings. Under
cross-family deployment, Qwen accuracy on the FFT+SigLIP2 disagreement subset ranges from
44.5\% to 74.1\% across the six pairings (Table~\ref{tab:cross_family_deployment}). In
several pairings, this is below FFT-only accuracy on the same target dataset, which
explains why escalation does not reliably beat the stronger cheap detector in that
setting. The verifier provides useful signal on some routed subsets, but its value is
distribution-dependent.

\begin{figure}[h]
    \centering
    \includegraphics[width=.95\columnwidth]{plot_fig_b_cascade_gain_scatter.png}
    \caption{Cascade gain over the stale semantic probe versus stale semantic accuracy on SDXL and FLUX. Each point corresponds to one semantic backbone and deployment dataset.}
    \label{fig:cascade_gain}
\end{figure}

Figure~\ref{fig:cascade_gain} shows the same pattern for GPT-5.5 under SDXL and FLUX
deployment shift. Cascade gain generally increases as the stale semantic probe becomes
less accurate. GPT-5.5 escalation gives positive gain in five of six configurations,
with FLUX using CLIP as the only negative-gain case. This motivates validating verifier
strength on the routed disagreement subset instead of enabling escalation from
disagreement rate alone.

For SigLIP2 and DINOv2, paired bootstrap confidence intervals over 10{,}000 resamples
confirm that three of four GPT-5.5 gains are significant at 95\%. The exception is FLUX
with SigLIP2, where the gain of $+$0.004 has a 95\% confidence interval of
[$-$0.004, $+$0.012], so it should be treated as a near-tie. CLIP gains are reported
without bootstrap intervals because per-sample stale predictions were not retained.

\subsection{Generator Distribution Shift}

CNNSpot accuracy is 54.0\%, 50.2\%, and 72.1\% across the three datasets. The
near-chance performance on GenBuster and SD~1.4 confirms that artifact-based detectors
trained on GANs do not generalize to modern diffusion and video
generators~\cite{Wang_2023_ICCV}. The FFT probe generalizes substantially better:
73.5\%, 74.8\%, and 99.8\%. Its spectral features capture regularities common across
synthesis architectures rather than GAN-specific upsampling patterns. The cascade
(FFT+SigLIP2, Qwen verifier) improves over FFT-only by 16.3 and 11.9 percentage
points on GenBuster and SD~1.4. On BigGAN, FFT-only already achieves 99.8\% and the
cascade matches it exactly.




\subsection{Limitations}

One validity concern is the real-image source used in the SDXL/FLUX deployment-shift
evaluation. The real half of these benchmarks uses ImageNet val/nature photographs
rather than the bundled AIGIBench PASCAL VOC images (Section~IV.B). This means the real
images come from the same source as the real images used to train the SD~1.4 semantic
probes, while the fake SDXL/FLUX images remain out of distribution.

Class-wise results suggest that this real-image alignment does not explain the reported
cascade gain. The stale semantic probe already performs near ceiling on the real class
(95.8\% on SDXL and 95.4\% on FLUX). Escalating disagreement cases to GPT-5.5 lowers
real-class accuracy by 4.0--4.3 points rather than increasing it. The accuracy recovery
comes from the fake class: stale fake-class accuracy rises from 73.6\% to 84.7\% on SDXL
and from 86.5\% to 91.2\% on FLUX. Thus, the reported gain comes from the
out-of-distribution half of the benchmark, while the real-image substitution works
against cascade accuracy.

The evaluation sets are modest, with 400 samples for GenBuster and 480 each for SD~1.4
and BigGAN. Bootstrap confidence intervals confirm statistical reliability at these
scales, but larger sets would strengthen generalization claims. The bandit router is
trained and evaluated within the same benchmark distribution, so its generalization
to unseen generator types remains untested. Qwen2.5-VL is applied zero-shot without
fine-tuning. However, a fine-tuned VLM may recover the domain sensitivity failures observed on
SD~1.4 and BigGAN. The study uses frame-level detection without temporal modeling,
and incorporating motion and frame consistency could improve performance on
video-specific generators. The semantic probe does not transfer across generator
families and deploying this cascade on a new family requires retraining the probe on
in-domain data. The disagreement rule itself has no trainable parameters and needs
no retraining.

\subsection{Practical Implications}

Disagreement-based routing provides a simple deployment signal for adaptive detection
pipelines. It requires no confidence calibration, learned policy, or retraining of the
routing rule itself. When a well-trained in-domain semantic probe is available, using the
semantic probe without VLM escalation gives the best cost-accuracy tradeoff. When the
semantic probe is stale under generator evolution, routing disagreement cases to a strong
verifier can recover accuracy while avoiding most verifier calls. On SDXL and FLUX,
GPT-5.5 escalation recovers up to 4.9 percentage points over the stale semantic probe
while escalating only 35--42\% of inputs (Table~\ref{tab:e8_results}).

The cascade is useful only when the verifier is stronger than the stronger cheap detector
on the routed disagreement subset. Under cross-family deployment
(Table~\ref{tab:cross_family_deployment}), FFT-only outperforms the full cascade in five
of six pairings, so a practitioner without an in-domain semantic probe should prefer the
strong cheap detector over escalating cross-family disagreements to Qwen. The BigGAN
distant-shift result (Section~V.D.7) reinforces this guideline that none of the tested
verifiers, including GPT-5.5, beats the stale semantic probe on the routed subset. This
contrasts with the SD~1.4-to-SDXL/FLUX setting, where GPT-5.5 provides positive gain
in five of six configurations. The practical decision therefore depends on the verifier
strength and shift type, rather than on disagreement alone.

\subsection{A Practical Fail-Safe: the Verifier Validation Gate}
\label{sec:validation_gate}

A deployed system may not know whether an incoming distribution is in-domain, a benign
generator update, or a larger cross-family shift. A natural first option is to use the
disagreement rate as a shift signal, escalating only when it exceeds the 26--28\%
in-domain baseline (Section~V.C). This is insufficient. Cross-family disagreement rates
range from 38.5--50.4\%, while SDXL/FLUX disagreement rates range from 35.2--41.4\%.
These ranges overlap even though escalation hurts in the cross-family setting and helps
with GPT-5.5 in most SDXL/FLUX configurations. BigGAN distant shift gives the opposite
failure mode: its disagreement rate is only 29.3\%, close to the in-domain baseline, even
though every tested verifier hurts performance there.

We therefore propose a \textbf{verifier validation gate}. Before allowing escalation for
a verifier $V$ on a new target distribution, evaluate $V$ on a small labeled calibration
sample drawn from the disagreement set. Enable the cascade only if
$A_{V,\text{cal}} > \max(A_{1,\text{cal}}, A_{2,\text{cal}})$ on that same calibration
sample. Otherwise, fall back to the stronger cheap detector. This rule does not require
knowing the generator family or shift type. It only requires labels for a small sample of
the routed subset.

The retrospective validation includes six cross-family pairings with Qwen2.5-VL-7B
(Table~\ref{tab:cross_family_deployment}), six SDXL/FLUX configurations under each of
GPT-5.5 and Qwen2.5-VL-7B (Table~\ref{tab:e8_results}), and three BigGAN distant-shift
configurations under GPT-5.5, Qwen2.5-VL-7B, and Qwen2.5-VL-72B
(Table~\ref{tab:biggan_shift}). Across these 21 configurations, using the full
disagreement subset as an oracle calibration sample makes the correct escalate-or-fallback
decision in 20 cases (95.2\%). The only miss is FLUX with SigLIP2 and GPT-5.5, a
near-tie where the verifier and semantic probe differ by 0.04 percentage points
(90.88\% vs.\ 90.93\%).

The gate correctly rejects escalation for the BigGAN distant shift and correctly enables
GPT-5.5 escalation in five of six SDXL/FLUX configurations, without using generator
identity or shift labels. In deployment, a smaller calibration sample would introduce
estimation noise. We report the oracle version as an upper bound on the rule's achievable
reliability and leave calibration-sample-size sensitivity to future work.



\section{Conclusion}

This work frames AI-generated media detection as an adaptive model selection problem. A two-stage cheap pipeline routes inputs through an FFT artifact detector and a semantic probe and escalates to a frontier VLM only when the two detectors disagree. Three research questions were evaluated with this design.

For RQ1, FFT and semantic probe agreement is a near-perfect predictor of detection ease. FFT accuracy is 99.7\%, 96.6\%, and 100\% on agreement cases and falls to 4.5\%, 17.4\%, and 66.7\% on disagreement cases. This pattern holds across video diffusion, image diffusion, and GAN generator families and forms the basis for the routing rule.

For RQ2, the disagreement rule does not require parameters or calibration. A bandit
router trained with 11 detector-derived features and an explicit latency reward, using
the same FFT+SigLIP2+Qwen backbone as the disagreement rule, matches the disagreement
rule's performance on SD~1.4 and BigGAN and is only slightly higher on GenBuster. This
suggests that the learned policy mainly recovers the same dominant routing signal rather
than discovering a consistently stronger one.

For RQ3, disagreement-based escalation reduces verifier cost, but its accuracy benefit is
conditional. On SDXL and FLUX, GPT-5.5 escalation on disagreement cases recovers up to
4.9 percentage points over a stale semantic probe while avoiding 58--65\% of verifier
invocations (Table~\ref{tab:e8_results}). Under cross-family transfer across
GenBuster/SD~1.4/BigGAN with Qwen as verifier, escalation underperforms FFT-only in five
of six pairings (Table~\ref{tab:cross_family_deployment}). Under the more distant BigGAN
shift, escalation is significantly harmful for all three tested verifiers, including
GPT-5.5 (Section~V.D.7). Thus, RQ3's answer is conditional: escalation helps only when
the verifier is stronger than the stronger cheap detector on the routed disagreement
subset.

The effective deployment of heterogeneous detection pipelines depends on the use of cheap detectors to identify when expensive verification is needed rather than increasing the complexity of the individual detector. Future work will explore temporal modeling for video detection, fine-tuned VLM verifiers for domain-specific forensic cues, and cross-family semantic probe training to remove the per-family retraining requirement.


\bibliographystyle{IEEEtran}
\bibliography{references}
\vskip -.5\baselineskip plus -1fil
\begin{IEEEbiography}[{\includegraphics[width=1in,height=1.25in,clip,keepaspectratio]{ansuman_nayak2.png}}]{Ansuman Nayak}
is a Software Engineer at Bank of America, where he focuses on building scalable automation tools for deployment pipelines, data workflows, and infrastructure operations. He earned his Bachelor's degree in Computer Science from Virginia Tech and is currently pursuing a Master's degree in Computer Science at Georgia Institute of Technology, specializing in Artificial Intelligence. His interests center on Agentic AI and LLM optimization, with a focus on workflow automation designed to operate efficiently under resource and cost constraints.
\end{IEEEbiography}
\vspace{-0.5\baselineskip}
\begin{IEEEbiography}[{\includegraphics[width=1in,height=1.25in,clip,keepaspectratio]{vijay.png}}]{Vijay K. Madisetti}
holds the position of Professor of Cybersecurity and Privacy (SCP) at Georgia Tech.
He earned his PhD in Electrical Engineering and Computer Sciences from the University of California at Berkeley
and is recognized as a Fellow of the IEEE. Additionally, he has been honored with the Terman Medal
by the American Society of Engineering Education (ASEE). Professor Madisetti has authored several widely
referenced textbooks on topics including cloud computing, data analytics, blockchain, and microservices.
\end{IEEEbiography}

\EOD

\end{document}
